\chapter{Elementary Reductions of the Riemann Hypothesis}
\label{v4-arith:w5-direct:closure-path}

The arithmetic Koszul reductions isolate the zero-placement problem
inside the Weil distribution.  The classical equivalences are recalled
with their hypotheses, and the restricted Paley--Wiener class is kept
separate from the full Weil cone.

\section{The Riemann Hypothesis in Classical Form}
\label{v4-arith:w5-direct:RH-classical}

\begin{definition}[the zeta function]
\label{v4-arith:w5-direct:def:zeta}
For \(\mathrm{Re}(s)>1\), the Riemann zeta function is the absolutely
convergent sum
\begin{equation}
\label{v4-arith:w5-direct:eq:zeta-dirichlet}
\zeta(s)\;:=\;\sum_{n=1}^{\infty}\frac{1}{n^{s}}.
\end{equation}
Its Euler product, equivalent on \(\mathrm{Re}(s)>1\), is
\begin{equation}
\label{v4-arith:w5-direct:eq:zeta-euler}
\zeta(s)\;=\;\prod_{p}\bigl(1-p^{-s}\bigr)^{-1},
\end{equation}
the product taken over all prime numbers.
\end{definition}

\begin{definition}[completed zeta and xi]
\label{v4-arith:w5-direct:def:xi}
Set
\begin{align}
\label{v4-arith:w5-direct:eq:Lambda}
\Lambda_{\zeta}(s)
&\;:=\;\pi^{-s/2}\,\Gamma(s/2)\,\zeta(s),
\\[2pt]
\label{v4-arith:w5-direct:eq:xi}
\xi_{\mathbb{R}}(s)
&\;:=\;\tfrac{1}{2}\,s(s-1)\,\Lambda_{\zeta}(s).
\end{align}
The function \(\Lambda_{\zeta}\) is meromorphic with simple poles at
\(s=0\) and \(s=1\); \(\xi_{\mathbb{R}}\) is entire. Both satisfy the
functional equation
\begin{equation}
\label{v4-arith:w5-direct:eq:functional-equation}
\Lambda_{\zeta}(s)\;=\;\Lambda_{\zeta}(1-s),
\qquad
\xi_{\mathbb{R}}(s)\;=\;\xi_{\mathbb{R}}(1-s).
\end{equation}
\end{definition}

\begin{remark}[functional equation via Poisson summation]
\label{v4-arith:w5-direct:rem:Poisson}
Let $\theta(t)=\sum_{n\in\mathbb Z}e^{-\pi n^2t}$ for $t>0$.
Poisson summation gives $\theta(1/t)=t^{1/2}\theta(t)$.
For $\Re s>1$, termwise integration gives
$\Lambda_\zeta(s)=\frac12\int_0^\infty(\theta(t)-1)t^{s/2}dt/t$.
Splitting at $1$ and substituting $t\mapsto1/t$ in the lower
integral yields its meromorphic continuation
\[
 \Lambda_\zeta(s)=\frac1{s(s-1)}+
 \frac12\int_1^\infty(\theta(t)-1)
       \bigl(t^{s/2}+t^{(1-s)/2}\bigr)\frac{dt}{t}.
\]
The remaining integral is entire by exponential decay. This formula
is invariant under $s\mapsto1-s$; multiplication by $s(s-1)/2$
removes the two poles and proves the functional equation for
$\xi_{\mathbb R}$.
Sources: B.~Riemann, ``\"{U}ber die Anzahl der Primzahlen unter einer gegebenen Gr\"{o}\ss e'', Monatsber.~Berlin.~Akad.~(1859); E.~C.~Titchmarsh, ``The Theory of the Riemann Zeta-Function'', Oxford 1951, Ch.~2.
\end{remark}

\begin{conjecture}[Riemann Hypothesis]
\label{v4-arith:w5-direct:conj:RH}
All non-trivial zeros of \(\zeta\) have real part \(1/2\). Equivalently,
the zeros of \(\xi_{\mathbb{R}}\) have real part \(1/2\).
\end{conjecture}

\section{Four Classical Equivalents}
\label{v4-arith:w5-direct:four-equivalents}

Four classical restatements of
\Cref{v4-arith:w5-direct:conj:RH} are used below.  They do not
identify the same object; they identify the same zero-placement
condition through different test categories.


\begin{theorem}[Weil positivity]
\label{v4-arith:w5-direct:thm:weil-positivity}
Let $\mathscr A=C_c^\infty(\mathbb R)$ with its inductive-limit
 topology, and put $f^*(t)=\overline{f(-t)}$ and
$h_f(z)=\int_{\mathbb R}f(t)e^{izt}dt$.
The Mellin coordinates are $\Phi_f(s)=h_f(-i(s-1/2))$.
With $\phi(x)=f(\log x)$, the change of variables $x=e^t$ gives
$\Phi_f(s)=\int_0^\infty\phi(x)x^{s-1/2}dx/x$.
For the zeros $\rho$ of $\xi_{\mathbb R}$, counted with multiplicity,
put $z_\rho=-i(\rho-1/2)$. Then
\begin{equation}
\label{v4-arith:w5-direct:eq:weil-distribution}
 W(f)=\sum_\rho h_f(z_\rho)
     =h_f(i/2)+h_f(-i/2)-W_\infty(f)-W_{\mathrm{fin}}(f),
\end{equation}
where
\[
\begin{aligned}
 W_\infty(f)&=\frac1{2\pi}\int_{\mathbb R}h_f(u)
 (\log\pi-\Re\psi(1/4+iu/2))du,\\
 W_{\mathrm{fin}}(f)&=\sum_{m\ge2}\frac{\Lambda(m)}{\sqrt m}
 [f(\log m)+f(-\log m)].
\end{aligned}
\]
Write $P(f)=h_f(i/2)+h_f(-i/2)$ for the polar functional.
Here $\Lambda(m)$ is the von Mangoldt function and $\psi=\Gamma'/\Gamma$.
The prime sum is finite. The zero sum converges absolutely.
RH is equivalent to $W(f*f^*)\ge0$ for every $f\in\mathscr A$.
\end{theorem}
\begin{proof}
These are the compact-test formula and Weil criterion of
\Cref{v4-arith:w5-a:thm:weil-explicit}, \Cref{v4-arith:w5-a:cor:weil-positivity-RH}.
The Fourier image of $f*f^*$ is $h_fh_f^\sharp$, where
$h^\sharp(z)=\overline{h(\bar z)}$. On the real line this product
is $|h_f|^2$.
\end{proof}
\begin{theorem}[Li positivity]
\label{v4-arith:w5-direct:thm:li-positivity}
\ClaimStatusProvedElsewhere. Define
\begin{equation}
\label{v4-arith:w5-direct:eq:li-coef}
\lambda_{n}\;:=\;\sum_{\rho}\left(1-\Bigl(1-\frac{1}{\rho}\Bigr)^{n}\right)
\;=\;\frac{1}{(n-1)!}\frac{d^{n}}{ds^{n}}
\bigl[s^{n-1}\log\xi_{\mathbb{R}}(s)\bigr]\Big|_{s=1}.
\end{equation}
Then RH holds if and only if \(\lambda_{n}\geq 0\) for all \(n\geq 1\).
Source: X.-J. Li, ``The positivity of a sequence of numbers and the
Riemann hypothesis'', J.~Number Theory 65 (1997), 325--333.
\end{theorem}

Let $\gamma$ denote the Euler--Mascheroni constant. For a normal
Hilbert--Schmidt operator $T$, its regularized determinant is
\[
 \det{}_2(I-sT)=\prod_j(1-s\lambda_j)e^{s\lambda_j},
\]
where the nonzero eigenvalues $\lambda_j$ are repeated with their
multiplicities. The product converges locally uniformly in $s\in\mathbb C$:
$\sum_j|\lambda_j|^2<\infty$ and
$(1-s\lambda_j)e^{s\lambda_j}=1+O_K(|\lambda_j|^2)$ uniformly on
each compact set $K\subset\mathbb C$.

\begin{theorem}[normalized Hilbert--P\'olya spectral form]
\label{v4-arith:w5-direct:thm:HP-form}
RH is equivalent to the existence of a separable Hilbert space
$H$ and a densely defined self-adjoint operator $A$ such that, for
$\Theta=\tfrac12I+iA$, the inverse $\Theta^{-1}$ is Hilbert--Schmidt
and the prescribed entire determinant is
\[
 \frac12e^{Bs}\det{}_2(1-s\Theta^{-1})=\xi_{\mathbb R}(s),
 \qquad B=-1-\frac\gamma2+\frac12\log(4\pi).
\]
Zeros of the determinant are counted with their spectral
multiplicities. The prescription fixes both its value $1/2$ at zero
and its logarithmic derivative $B$ there.
\end{theorem}
\begin{proof}
For a self-adjoint $A$, the bounded inverse of $\tfrac12I+iA$
is normal. If it is Hilbert--Schmidt it is compact; its nonzero
eigenvalues are $(1/2+ia_j)^{-1}$ for real eigenvalues $a_j$,
with their multiplicities. The regularized product therefore has
zeros only at $1/2+ia_j$. Its asserted equality with $\xi_{\mathbb R}$
implies RH.
Conversely, on the zero-occurrence carrier $K_Z=\ell^2(Z)$ of
\Cref{v4-arith:sc:determinant}, use the closed normal operator
$\Theta_Z$ with maximal domain
$\{u:\sum_\rho|\rho u_\rho|^2<\infty\}$.
Under RH write $\rho=1/2+i\gamma_\rho$, with $\gamma_\rho\in\mathbb R$.
The diagonal operator
\[
 (Au)_\rho=\gamma_\rho u_\rho,\qquad
 \operatorname{Dom}A=\{u:\sum_\rho|\gamma_\rho u_\rho|^2<\infty\}
\]
is self-adjoint, and $\Theta_Z=\tfrac12I+iA$ on the same domain.
Its inverse is Hilbert--Schmidt, and
\Cref{v4-arith:sc:normalized-det} proves exactly the stated
normalized genus-one determinant. That construction uses the zero
divisor as input. Constructing an arithmetic operator and an
intertwiner with this carrier remains an additional problem.
\end{proof}
The Hilbert--P\'olya question asks for such an independent arithmetic
realization; the abstract zero-occurrence construction does not supply
it. Connes, \emph{Selecta Mathematica} 5 (1999), arXiv:math/9811068,
treats a separate trace-formula framework.

\paragraph{Entire functions and the zero divisor.}
Write $F(z)=\Xi(z)=\xi_{\mathbb R}(1/2+iz)$ and
$G^\#(z)=\overline{G(\bar z)}$ for an entire function $G$.
The functional equation and reality of $\xi_{\mathbb R}$ make $F$
even and real entire. RH says precisely that all its zeros are real.
Use the convention that an entire function $E$ belongs to
$\mathrm{HB}$ if it has no real zeros and
$|E^\#(z)|<|E(z)|$ for $\Im z>0$.
No finite-exponential-type condition is included in this definition.

\begin{theorem}[a shift-family Hermite--Biehler criterion]
\label{v4-arith:w5-direct:thm:dB-form}
For $h>0$ put
\[
 E_h(z)=F(z+ih)=\xi_{\mathbb R}(1/2+h-iz).
\]
The following conditions are equivalent: RH; $E_h\in\mathrm{HB}$
for every $h>0$; and $E_{h_n}\in\mathrm{HB}$ for some positive
sequence $h_n\downarrow0$. The recovery identities are
\[
 E_h^\#(z)=F(z-ih),\qquad F(z)=E_h(z-ih).
\]
Thus a zero $w$ of $E_h$ gives the zero $w+ih$ of $F$ with exactly
the same multiplicity. No simplicity assumption is needed.
\end{theorem}
\begin{proof}
One has $F(0)\ne0$. Indeed the alternating series at $1/2$ can be
paired as
$\eta(1/2)=\sum_{n\ge1}[(2n-1)^{-1/2}-(2n)^{-1/2}]>0$;
$\zeta(1/2)=\eta(1/2)/(1-\sqrt2)$ is consequently nonzero.
Under RH the genus-one factorization and evenness give
\[
 F(z)=F(0)\prod_{\gamma>0}
       (1-z^2/\gamma^2)^{m_\gamma},\qquad
 \sum_{\gamma>0}m_\gamma\gamma^{-2}<\infty.
\]
Here $\pm\gamma$ are the zeros of $F$, with multiplicity
$m_\gamma$. The inverse-square convergence follows from
\Cref{v4-arith:pm:zero-count}. Pairing the factors at $\gamma$ and
$-\gamma$ cancels their genus-one exponentials; the remaining linear
exponential is constant by evenness. The product converges locally
uniformly.
For $z=x+iy$, $y>0$, every real $\gamma$ satisfies
\[
 |z+ih-\gamma|^2-|z-ih-\gamma|^2=4yh>0.
\]
Apply this to both members of every paired factor. If
$F(z-ih)\ne0$, each factor in the ratio of absolute values exceeds
one. Its limiting product remains strictly greater than one by
retaining one strict factor and bounding all other factors below by
one. If $F(z-ih)=0$, the inequality is immediate, since $F(z+ih)$
cannot vanish in the upper half-plane under RH.
Hence $|E_h(z)|>|E_h^\#(z)|$.
All zeros of $E_h$ are $\gamma-ih$, so none is real.
This proves the assertion for every $h>0$, and therefore for any
sequence decreasing to zero.
Conversely, a nonreal zero of the real entire function $F$ gives a
zero $w$ in the upper half-plane. For $h_n<\Im w$, the point
$w-ih_n$ is an upper-half-plane zero of $E_{h_n}$, contradicting the
strict HB inequality. The two recovery identities follow directly
from reality, translation and the functional equation.
\end{proof}

\begin{remark}[one shift and a family of shifts]
\label{v4-arith:consumer53:fixed-shift}
For $F_a(z)=z^2+a^2$, $a>0$, and $h>a$,
\[
 F_a(z+ih)=(z+i(h-a))(z+i(h+a))\in\mathrm{HB},
\]
although $F_a$ has nonreal zeros. Each factor has the strict
upper-half-plane modulus inequality because
$|z+ib|^2-|z-ib|^2=4b\Im z>0$ for $b>0$.
Thus one fixed shift does not characterize real-rootedness.
For $\Xi$ itself, Lagarias, arXiv:math/0601653, Lemma~2.1,
pp.~5--6, proves the shift inequality unconditionally for
$h\ge1/2$, and under RH for every $h>0$; the change of variable is
$s=1/2-iz$. Section~6 relates these functions to de Branges spaces.
The family tending to zero, not one large shift, is the RH criterion.
\end{remark}

\begin{theorem}[the derivative criterion and real multiplicities]
\label{v4-arith:consumer53:derivative-HB}
Set $D(z)=F(z)+iF'(z)$. Then
\[
 D^\#(z)=F(z)-iF'(z),\qquad
 F(z)=\tfrac12(D(z)+D^\#(z)).
\]
RH is equivalent to $|D(z)|>|D^\#(z)|$ for every $\Im z>0$,
where real zeros of $D$ are allowed. With the no-real-zero convention
above, $D\in\mathrm{HB}$ is instead equivalent to RH together with
simplicity of all zeros of $F$.
\end{theorem}
\begin{proof}
Under RH, logarithmic differentiation of the paired product gives
for $z=x+iy$, $y>0$,
\[
 \Im\frac{F'(z)}{F(z)}
 =-y\sum_{\gamma\in\mathbb R}
           \frac{m_\gamma}{|z-\gamma|^2}<0.
\]
The imaginary-part series is absolutely convergent by the
inverse-square estimate. There is at least one zero, indeed
infinitely many by \Cref{v4-arith:pm:zero-count}. Consequently
\[
 |D(z)|^2-|D^\#(z)|^2
 =-4|F(z)|^2\Im\frac{F'(z)}{F(z)}>0.
\]
Conversely, at an upper-half-plane zero $w$ of $F$ the two values
$D(w)=iF'(w)$ and $D^\#(w)=-iF'(w)$ have equal modulus.
This contradicts the strict inequality, also when both vanish at a
multiple zero. Conjugation symmetry excludes lower-half-plane zeros
as well.
For a real zero write $F(z)=(z-\gamma)^m g(z)$ with $g(\gamma)\ne0$.
Then
\[
 D(z)=(z-\gamma)^{m-1}
 \bigl((z-\gamma)g(z)+i[m g(z)+(z-\gamma)g'(z)]\bigr).
\]
The bracket is nonzero at $\gamma$, so the real multiplicity in
$D$ is exactly $m-1$. On the real axis $F,F'$ are real, hence a
real zero of $D$ occurs only where both vanish. This proves the
simplicity qualification and the stated recovery identity.
\end{proof}

\begin{proposition}[the growth required by the zeta divisor]
\label{v4-arith:consumer53:exponential-type}
No nonzero entire function of finite exponential type can have a
zero divisor containing the full divisor of $F=\Xi$.
In particular no nonzero real or imaginary component of a
finite-exponential-type function can have that divisor.
The functions $F$ and each fixed translate $E_h$ have order one
and infinite exponential type.
\end{proposition}
\begin{proof}
Let $N(T)$ count the nontrivial zeta zeros with $0<\Im\rho\le T$,
with multiplicity. The Riemann--von Mangoldt estimate is
$N(T)=\frac{T}{2\pi}\log\frac{T}{2\pi e}+O(\log T)$.
For a precise primary bound see Hasanalizade--Shen--Wong,
arXiv:2107.06506, Corollary~1.2, equation~(1.5), p.~2.
The strip $0\le\Re\rho\le1$ and conjugation give, for the number
$n_F(r)$ of zeros in $|z|\le r$,
\[
 2N(\sqrt{r^2-1/4})\le n_F(r)\le2N(r),\qquad
 n_F(r)=\frac r\pi\log\frac r{2\pi e}+O(\log r).
\]
If $G$ is nonzero entire with $|G(z)|\le Ce^{\tau|z|}$ and
$G(0)\ne0$, Jensen's formula at radius $2r$ gives
\[
 n_G(r)\log2\le
 \frac1{2\pi}\int_0^{2\pi}\log|G(2re^{i\theta})|d\theta
 -\log|G(0)|=O(r).
\]
Factoring out a zero at the origin changes the count by a constant.
Thus $G$ cannot contain the $r\log r$ divisor of $F$.
If $E$ has finite exponential type, so do its components
$(E+E^\#)/2$ and $(E^\#-E)/(2i)$; apply the same argument to every
nonzero component.
Finally the functional equation, Stirling's formula and
$\zeta(\sigma)\to1$ give
\[
 \log F(iy)=\log\xi_{\mathbb R}(1/2+y)
 =\frac y2\log y-\frac y2(1+\log(2\pi))+O(\log y)
 \quad(y\to+\infty).
\]
Together with the order-at-most-one bound in
\Cref{v4-arith:pm:zero-count}, this proves order one and infinite
exponential type, also after any fixed translation.
\end{proof}
The general Hilbert-space theory of entire functions is developed in
de Branges, \emph{Hilbert Spaces of Entire Functions} (1968).
The explicit constructions above specify the recovery of the zeta
divisor and retain its actual growth and multiplicities.

\begin{remark}[four precise equivalences]
\label{v4-arith:w5-direct:rem:four-equivalents}
Weil positivity on all compact tests, Li positivity, the normalized
spectral realization of \Cref{v4-arith:w5-direct:thm:HP-form}, and
the shift-family condition of \Cref{v4-arith:w5-direct:thm:dB-form}
are each equivalent to RH. The derivative inequality gives the
same zero condition with its real-zero convention stated explicitly.
These criteria use different carriers: test functions, a coefficient
sequence, a spectral operator and entire functions. None of their
abstract equivalences constructs the native arithmetic pairing or
operator maps required below.
\end{remark}

\section{The Arithmetic Reduction}
\label{v4-arith:w5-direct:programme-elementary}

The arithmetic Koszul construction supplies five comparison objects.
Each object is defined on its own domain, and no implication from a
restricted domain to the full Weil cone is used without an extension
theorem.

\subsection{Step 1: The arithmetic Heisenberg Koszul pair}
\label{v4-arith:w5-direct:step-1}

\begin{question}[the arithmetic Heisenberg coefficient object]
\label{v4-arith:w5-direct:def:heisenberg}
Construct local coefficient objects \(\mathsf H_v\), a common
coefficient category, pointings and a completion that define
\begin{equation}
\label{v4-arith:w5-direct:eq:heis}
 \mathsf H^{\mathrm{Ar}}
 =\mathsf H_\infty\widehat\otimes
                    \widehat\bigotimes_p\mathsf H_p.
\end{equation}
Compare the infinite-place factor with the Heisenberg field and
the finite factors with their continuous Galois coefficient
models. An associative bar construction additionally requires
an augmented algebra, or a specified curved replacement.
The pointed tensor construction of
\Cref{v4-arith:tate:setup:data} applies once its inputs are given.
\end{question}

\begin{definition}[central invariants and comparison functors]
\label{v4-arith:w5-direct:def:koszul}
For a specified arithmetic object and its specified dual,
write \(\kappa^{\mathrm{Ar}}_{\mathsf H}\) and
\((\kappa^!)^{\mathrm{Ar}}_{\mathsf H}\) for their chosen
scalar central invariants. These are not the bar and cobar
complexes themselves.
An integral co--contra adjunction and an algebraic bar--cobar
adjunction have the distinct carriers described in
Chapter~\ref{v4-arith:arith-positselski}.
\end{definition}

\begin{theorem}[P1 under the central-defect comparison]
\label{v4-arith:w5-direct:thm:P1-elementary}
Assume the objects, involution, central-defect identity and
Mellin evaluation of Theorem~\ref{v4-arith:thm:P1-resolution}
realize these arithmetic invariants, with \(E(w)\ne0\).
Then
\[
 \kappa^{\mathrm{Ar}}_{\mathsf H}
       +(\kappa^!)^{\mathrm{Ar}}_{\mathsf H}=0
 \quad\Longleftrightarrow\quad
 \xi_{\mathbb R}(s)=\xi_{\mathbb R}(1-s).
\]
\end{theorem}
\begin{proof}
Applying \(E\) to the defect identity identifies the difference
of the two completed Mellin functions with
\((\kappa^{\mathrm{Ar}}_{\mathsf H}
 +(\kappa^!)^{\mathrm{Ar}}_{\mathsf H})E(w)\).
The nonzero meromorphic function \(E(w)\) detects the scalar
coefficient. Multiplication by the reflection-invariant
factor \(s(s-1)/2\) gives the equivalent identity for \(\xi\).
\end{proof}

\begin{remark}[the scalar and categorical comparisons]
\label{v4-arith:w5-direct:rem:P1-elementary}
Poisson summation proves the scalar functional equation.
An arithmetic interpretation requires maps identifying the
duality involution with Mellin reflection and the central
classes with the specified defect.
Equality of two scalar characters need not detect a nonzero
vector defect, as in Theorem~\ref{v4-arith:thm:P1-resolution}.
Neither the scalar equality nor a categorical equivalence
alone supplies those identifications.
\end{remark}

\subsection{Archimedean Heisenberg Character}
\label{v4-arith:w5-direct:step-2}

\begin{definition}[archimedean Heisenberg character]
\label{v4-arith:w5-direct:def:zhu}
The archimedean fibre \(\mathsf{H}_{\infty}\) is the standard Heisenberg
vertex operator algebra on \(\mathbb{C}_{\infty}\) with Virasoro zero-mode
\(L_{0}\) and central charge \(c_{\infty}=1\) (one free boson). For the
primary \(\mathsf{H}_{\infty}\)-module \(V_{\infty}\) (the vacuum
representation), the Zhu character is
\begin{equation}
\label{v4-arith:w5-direct:eq:zhu-char}
\chi_{V_{\infty}}(\tau)\;:=\;\mathrm{tr}_{V_{\infty}}\bigl(q^{L_{0}-c_{\infty}/24}\bigr),
\qquad q=e^{2\pi i\tau},
\end{equation}
a standard VOA Zhu trace on the archimedean fibre. The global object
\(\mathcal{V}^{\mathrm{prim}}\) is the Hochschild trace class of the
identity on the arithmetic Hecke category; it is \emph{not} a vertex
algebra. The quantity \(\Lambda_{\zeta}(s)\) arises as its Hochschild
character via Tate's adelic construction, not as a VOA module trace.
\end{definition}

\begin{theorem}[P2: Hochschild character and functional equation]
\label{v4-arith:w5-direct:thm:P2-elementary}
\ClaimStatusProvedHere. The Hochschild character of
\(\mathcal{V}^{\mathrm{prim}}\) equals \(\Lambda_{\zeta}(s)\) under the
identification of the Hochschild trace with Tate's adelic zeta integral
(\Cref{v4-arith:thm:G-row-Riemann-functional-equation}). The primary
determinant carries rank residue \(r_{\zeta}=2\), the residue at
\(s=1\) of \(2\zeta(s)\) (\Cref{v4-arith:thm:c-as-residue}). The adelic
functional equation \(\Lambda_{\zeta}(s)=\Lambda_{\zeta}(1-s)\) follows
from Poisson summation applied to the Schwartz--Bruhat function at each
place.  The Heisenberg vacuum character is the local eta-factor; the
theta transformation enters only after the lattice Schwartz sector is
adjoined.
\end{theorem}

\begin{remark}[elementary content of P2]
\label{v4-arith:w5-direct:rem:P2-elementary}
On the archimedean fibre,
\[
\chi_{V_{\infty}}(\tau)=q^{-1/24}\prod_{n\geq1}(1-q^{n})^{-1}.
\]
This eta-character is not Tate's theta integral.  Tate's local factor is
obtained from the Schwartz function on \(\mathbb{R}\) and the theta sum
\(\theta(t)=\sum_{n\in\mathbb{Z}}e^{-\pi n^{2}t}\); Poisson summation is
the identity \(\theta(t)=t^{-1/2}\theta(1/t)\).  No finite Zhu
\(S\)-matrix or rational-VOA modularity is used in the Hochschild trace.
Tate (1950) identified \(\Lambda_{\zeta}\) with the adelic zeta integral;
the present step repackages Tate's computation as a Hochschild-trace
identity on \(\mathcal{V}^{\mathrm{prim}}\), with VOA language confined
to \(\mathsf{H}_{\infty}\). Source: J.~T.~Tate, thesis, Princeton 1950
(reprinted in Cassels--Fr\"{o}hlich, ``Algebraic Number Theory'',
Academic Press 1967).
\end{remark}

\subsection{Step 3: Petersson positivity}
\label{v4-arith:w5-direct:step-3}


\begin{definition}[adelic Petersson form]
\label{v4-arith:w5-direct:def:petersson}
Fix a unitary central character and quotient Tamagawa measure $dh$ on
$[GL_2]=Z(\mathbb A)GL_2(\mathbb Q)\backslash GL_2(\mathbb A)$.
On the corresponding cuspidal $L^2$ space set
\begin{equation}
\label{v4-arith:w5-direct:eq:petersson}
 \langle v,w\rangle_{\mathrm{Pet}}=
 \int_{[GL_2]}v(h)\overline{w(h)}dh.
\end{equation}
For a proposed complex realization $E$ of the primary Hochschild
trace sector, a Petersson pairing on $E$ requires a complex-linear map
$I:E\to L^2_{\mathrm{cusp}}([GL_2],dh)$.
\end{definition}

\begin{theorem}[positive pullback and Koszul comparison]
\label{v4-arith:w5-direct:thm:P3-elementary}
The pullback $\langle Iv,Iw\rangle_{\mathrm{Pet}}$ is positive
semidefinite, and is definite precisely when $I$ is injective.
If $B_{\mathrm{Kos}}:E\times E\to\mathbb C$ is bilinear and
$\sigma:E\to E$ is an antilinear involution, the identity
\[
 B_{\mathrm{Kos}}(v,\sigma w)=\langle Iv,Iw\rangle_{\mathrm{Pet}}
\]
transports this positivity to the specified Koszul pairing.
\end{theorem}
\begin{proof}
Its diagonal is $\|Iv\|_2^2$, which vanishes exactly on $\ker I$.
The displayed equality then gives the last assertion.
\end{proof}

\begin{remark}[local and global realization data]
\label{v4-arith:w5-direct:rem:P3-elementary}
For a restricted Hilbert tensor product with unit reference vectors
$e_p$, the inner product on elementary tensors is
\begin{equation}
\label{v4-arith:w5-direct:eq:petersson-factorized}
 \langle v,w\rangle=\langle v_\infty,w_\infty\rangle
 \prod_p\langle v_p,w_p\rangle,
\end{equation}
where all but finitely many factors equal one.
This defines a positive Hilbert completion. Identifying it with
$E$ and with the adelic Petersson space requires local maps,
normalizations on $e_p$, and the global pairing identity above.
The primary Hochschild trace class is not itself a space of
functions on $[GL_2]$. Its degree, completion, realization map,
and pairing are the additional data specified in
\Cref{v4-arith:pres57:native-pairing}.
\end{remark}
\subsection{Step 4: Bost--Connes and the Hilbert--P\'olya Boundary}
\label{v4-arith:w5-direct:step-4}

\begin{definition}[Bost--Connes system and KMS convention]
\label{v4-arith:w5-direct:def:BC}
Let $D=C^*(\mathbb Q/\mathbb Z)$, with group unitaries $e(r)$,
and put
\[
 \alpha_n(e(r))=\frac1n\sum_{nt=r}e(t),\qquad n\in\mathbb N^\times.
\]
The Bost--Connes algebra is the semigroup crossed product
$\mathcal A_{\mathrm{BC}}=D\rtimes_\alpha\mathbb N^\times$.
Its multiplicative generators $\mu_n$ are isometries, with
\[
 \mu_n^*\mu_n=1,\quad \mu_m\mu_n=\mu_{mn},\quad
 \mu_n d\mu_n^*=\alpha_n(d),\quad
 \mu_n^*e(r)\mu_n=e(nr).
\]
The time evolution fixes $D$ and satisfies
$\sigma_t(\mu_n)=n^{it}\mu_n$.
For $\beta>0$ and entire analytic elements the KMS$_\beta$ identity is
$\varphi(ab)=\varphi(b\sigma_{i\beta}(a))$.
In particular
$\varphi(\mu_n\mu_n^*)=n^{-\beta}$ whereas
$\varphi(\mu_n^*\mu_n)=1$; for $n>1$ the KMS state is nontracial
and $\mu_n$ is not unitary.
The crossed-product description is that of Laca--Raeburn,
arXiv:math/9911134, \S2, p.~3.
\end{definition}

At $\beta=1$ let $\varphi_1$ be the unique KMS state and
$(\pi,H,\Omega)$ its GNS representation. Its restriction to $D$ is
$\varphi_1(e(r))=\delta_{r,0}$; see Connes,
\emph{Noncommutative Geometry}, V~\S11, Theorem~5(a), p.~528.
Invariance and the KMS identity give, on the dense monomial algebra,
\[
 \varphi_1(\mu_a e(r)\mu_b^*)
       =\delta_{a,b}\,a^{-1}\delta_{r,0}.
\]
Indeed unequal $a,b$ give a nontrivial time character, while for
$a=b$ one moves $\mu_a$ around the state by the KMS identity.
For this canonical state its normal extension to
$M=\pi(\mathcal A_{\mathrm{BC}})''$ is faithful, and $\Omega$ is
cyclic and separating. We use its Tomita modular operator $\Delta_1$,
with support $1$ in $M$ and convention
\[
 \Delta_1^{it}\pi(a)\Delta_1^{-it}=\pi(\sigma_{-t}(a)),
 \qquad \Delta_1^{it}\Omega=\Omega.
\]
This sign is the canonical Hecke-state convention in Connes,
V~\S11, Proposition~4, p.~527. In vectors below, $a\Omega$ means
$\pi(a)\Omega$.

\begin{theorem}[the full critical-temperature modular generator]
\label{v4-arith:w5-direct:thm:HPAr-elementary}
Define $N_{\mathrm{BC}}=-\tfrac12\log\Delta_1$ by self-adjoint
functional calculus. There is an orthogonal Hilbert sum
\[
 H=\widehat{\bigoplus}_{q\in\mathbb Q_{>0}}H_q,
 \qquad N_{\mathrm{BC}}|_{H_q}=\tfrac12\log q,
\]
where every $H_q$ has countably infinite dimension. Its maximal domain is
\[
 \operatorname{Dom}N_{\mathrm{BC}}
 =\left\{\xi=\sum_q\xi_q:
            \sum_q(\log q)^2\|\xi_q\|^2<\infty\right\}.
\]
Its point spectrum is $\tfrac12\log\mathbb Q_{>0}$ and its closed
spectrum is $\mathbb R$. All spectral measures are pure point.
\end{theorem}
\begin{proof}
Each monomial $x=\mu_a e(r)\mu_b^*$ is entire analytic with
$\sigma_t(x)=(a/b)^{it}x$. The modular identity therefore gives
$\Delta_1^{it}x\Omega=(a/b)^{-it}x\Omega$.
Take $H_q$ to be the closed span of these eigenvectors with $a/b=q$.
Different eigenvalues are orthogonal, and the monomials are dense
in the GNS space. Functional calculus gives the displayed direct
sum and maximal domain. Every $q=a/b$ occurs, since
\[
 \|\mu_a\mu_b^*\Omega\|^2
       =\varphi_1(\mu_b\mu_b^*)=b^{-1}>0.
\]
The closed spectrum is the closure of these diagonal eigenvalues,
which is $\mathbb R$ because the positive rationals are dense.
The diagonal spectral projections make every spectral measure
atomic on that countable set.
For multiplicity fix a prime $p$, and set
$P_k=\mu_{p^k}\mu_{p^k}^*$ and $d_k=P_k-P_{k+1}$, $k\ge0$.
The $P_k$ are nested fixed-point projections; hence the $d_k$ are
orthogonal projections with
$\|d_k\Omega\|^2=p^{-k}-p^{-k-1}>0$.
For fixed-point elements $x,y$, cancellation of $\mu_a^*\mu_a$ and
the KMS identity give
\[
 \varphi_1\bigl((\mu_a y\mu_b^*)^*(\mu_a x\mu_b^*)\bigr)
     =b^{-1}\varphi_1(y^*x).
\]
Thus $\mu_a d_k\mu_b^*\Omega$ are infinitely many nonzero orthogonal
vectors in $H_{a/b}$. The algebra, and hence $H$, is separable,
so its eigenspace dimensions are countably infinite.
\end{proof}

\begin{proposition}[a reducing occupation subspace]
\label{v4-arith:consumer53:occupation-subspace}
For an integer $s\ge1$ put
\[
 K_s=\overline{\operatorname{span}}
       \{\mu_{m^s}\Omega:m\in\mathbb N^\times\}.
\]
The displayed vectors are an orthonormal basis and $K_s$ reduces
$N_{\mathrm{BC}}$. The unitary
$V_s:\ell^2(\mathbb N^\times)\to K_s$,
$V_s\delta_m=\mu_{m^s}\Omega$, satisfies
\[
 V_s^*(N_{\mathrm{BC}}|_{K_s})V_s=\tfrac s2 H_{\mathrm{occ}},
 \qquad (H_{\mathrm{occ}}c)_m=(\log m)c_m,
\]
including the domain identity, where
\[
 \operatorname{Dom}H_{\mathrm{occ}}
       =\{c:\sum_m(\log m)^2|c_m|^2<\infty\}.
\]
Thus $K_1$ gives half the occupation energy and $K_2$ gives exactly
$\log m$.
\end{proposition}
\begin{proof}
Each vector has norm one because $\mu_{m^s}$ is an isometry.
For $m\ne n$, invariance of $\varphi_1$ gives
$\varphi_1(\mu_{m^s}^*\mu_{n^s})=0$, since $(n/m)^{ist}$ is a nontrivial time character. These orthonormal eigenvectors have energies
$(s/2)\log m$. Their closed span reduces every spectral projection
of $N_{\mathrm{BC}}$, proving reduction and the maximal diagonal
domain formula.
\end{proof}

\begin{proposition}[compressed shifts and the observable obstruction]
\label{v4-arith:consumer53:BC-compression}
Let $P_s$ be the orthogonal projection onto $K_s$, and set
$T_n=\pi(\mu_{n^s})|_{K_s}$. Under $V_s$, these operators and their
Hilbert adjoints act as
\[
 V_s^*T_nV_s\delta_m=\delta_{nm},\qquad
 V_s^*T_n^*V_s\delta_m=
 \begin{cases}\delta_{m/n},&n\mid m,\\0,&n\nmid m.\end{cases}
\]
The adjoint here is $T_n^*=P_s\pi(\mu_{n^s}^*)|_{K_s}$.
The original adjoints do not preserve $K_s$, and compression of the
BC algebra to $K_s$ is not a $*$-representation.
\end{proposition}
\begin{proof}
The semigroup identity gives
$\mu_{n^s}\mu_{m^s}\Omega=\mu_{(nm)^s}\Omega$.
Taking inner products of these basis vectors gives the adjoint formula.
For $n>1$ the vector $\mu_{n^s}^*\Omega$ has norm squared $n^{-s}$
and strictly negative modular energy, so it is nonzero and orthogonal
to $K_s$. In particular it cannot be the zero compressed adjoint
$T_n^*\Omega$.
For the group unitaries, time invariance makes the off-diagonal
matrix coefficients of $e(r)$ between the $K_s$ basis vectors vanish.
The transfer relation and the state on $D$ give the diagonal entries:
\[
 P_s\pi(e(r))|_{K_s}\,\mu_{m^s}\Omega
   =\mathbf1_{\{m^s r=0\ \mathrm{in}\ \mathbb Q/\mathbb Z\}}
       \mu_{m^s}\Omega.
\]
At $r=1/2$ the compressed operator kills $\Omega$, although $e(r)$
is unitary. Thus the unital compression preserves adjoints but does not preserve
products, and is not a representation of the BC algebra.
\end{proof}
Writing $m=\prod_p p^{k_p}$ identifies the occupation basis with
finitely supported nonnegative prime occupations. On their finite
linear span,
\[
 N_{\mathrm{BC}}|_{K_s}=\frac s2\sum_p(\log p)K_p,
 \qquad K_p\mu_{m^s}\Omega=k_p\mu_{m^s}\Omega.
\]
The closure is the maximal diagonal operator already specified.
This identifies a Hilbert space and its occupation generator; it
does not turn the compressed BC observables into a BC representation
or into bosonic creation operators.

\begin{remark}[occupation trace and the Hilbert--P\'olya boundary]
\label{v4-arith:w5-direct:rem:HPAr-stop}
On the reducing subspace $K_2$ the spectral theorem gives
\[
 \operatorname{Tr}_{K_2}(e^{-tN_{\mathrm{BC}}})
       =\sum_{m\ge1}m^{-t}=\zeta(t),\qquad t>1.
\]
At $t=1$ this trace diverges. On the full GNS space the exponential
$e^{-tN_{\mathrm{BC}}}$ is unbounded for every $t>0$, since its
spectrum contains all negative rational logarithmic energies.
There is also a separate standard BC representation on
$\ell^2(\mathbb N^\times)$, with
$\pi_0(\mu_n)\delta_m=\delta_{nm}$ and
$\pi_0(e(r))\delta_m=e^{2\pi i m r}\delta_m$.
The root-of-unity average in $\alpha_n$ is zero unless $n\mid m$,
and otherwise is $e^{2\pi i(m/n)r}$, proving the covariance relation.
The occupation Hamiltonian implements $\sigma_t$ there and has
partition function $\operatorname{Tr}(e^{-tH_{\mathrm{occ}}})=\zeta(t)$
for $t>1$. This is the type-I Gibbs realization underlying
Connes, V~\S11, Theorem~5(b), p.~528; it is not the compression of
$\pi$ to $K_2$, whose $e(r)$ matrices were computed above.
The occupation trace is consequently not the full modular trace
and not the signed Weil functional. A comparison with
\Cref{v4-arith:w5-direct:thm:HP-form}, including the arithmetic
question of \Cref{v4-arith:w5-b:hyp:twist-existence}, must supply
an actual test-algebra map, its involution and complete
$P-W_\infty-W_{\mathrm{fin}}$ pairing, an operator-domain comparison,
and the normalized zero divisor with its multiplicities.
The reducing subspace and its zeta trace do not supply those maps.
\end{remark}

\subsection{Restricted Weil pairings}
\label{v4-arith:w5-direct:step-5}

\begin{definition}[compact coefficient multiplier]
\label{v4-arith:w5-direct:def:mellin-lift}
For $f\in\mathscr A$, define
\begin{equation}
\label{v4-arith:w5-direct:eq:mellin-lift}
 U(f)(z)=\sum_{n\ge1}f(\log n)n^{1/2}e^{2\pi inz}.
\end{equation}
This is a finite holomorphic Fourier polynomial with exponential
decay as $\Im z\to\infty$.
\end{definition}

\begin{proposition}[modularity of the compact multiplier]
A nonzero $U(f)$ cannot be a cusp form of positive weight for
$SL_2(\mathbb Z)$.
\end{proposition}
\begin{proof}
A finite Fourier polynomial is analytic in $y$ at $z=iy=0$.
The modular transformation $z\mapsto-1/z$ would imply
$U(f)(iy)=O(y^{-k}e^{-2\pi/y})$. All its Taylor coefficients
would vanish. This is the compact multiplier obstruction of
\Cref{v4-arith:w5-a:lem:upsilon-cuspidal}.
\end{proof}

\begin{theorem}[restricted positive realizations]
\label{v4-arith:w5-direct:thm:WP-isometry}
Let $D\subset\mathscr A$ be a complex vector subspace.
Put $q_D(f,g)=W(f*g^*)$ for $f,g\in D$.
A linear map $J:D\to K$ to a positive Hilbert space satisfying
\begin{equation}
\label{v4-arith:w5-direct:eq:WP-isometry}
 \langle Jf,Jg\rangle_K=q_D(f,g)
 =P(f*g^*)-W_\infty(f*g^*)-W_{\mathrm{fin}}(f*g^*)
\end{equation}
exists exactly when $q_D$ is positive semidefinite.
When it exists, the completion of $D/\{f:q_D(f,f)=0\}$ gives
its minimal realization.
\end{theorem}
\begin{proof}
A Hilbert realization gives positivity. Conversely the
Cauchy--Schwarz inequality for a positive semidefinite Hermitian
form makes its nullspace orthogonal to $D$. The quotient is an
inner-product space, and its completion provides $J$.
\end{proof}

\begin{corollary}[restricted positivity from a native pairing]
\label{v4-arith:w5-direct:thm:HPLi-elementary}
Suppose $T:D\to E$ and $I:E\to L^2_{\mathrm{cusp}}([GL_2],dh)$
are linear maps with
$W(f*g^*)=\langle ITf,ITg\rangle_{\mathrm{Pet}}$ for $f,g\in D$.
Then $W(f*f^*)\ge0$ on $D$.
\end{corollary}
\begin{proof}
Take $g=f$. The resulting value is $\|ITf\|_2^2$.
\end{proof}
The construction of $T,I$ and this equality is the restricted
arithmetic comparison problem. Neither the multiplier $U$ nor
moment conditions supply these maps. A restricted Li trace additionally
requires explicit Li test approximants and convergence of their zero
sums. The rational Li tests are not compact Fourier tests, by
\Cref{v4-arith:w5-a:prop:li-via-weil}.

\section{The restricted test class and its closure}
\label{v4-arith:w5-direct:closing-restricted}
\begin{theorem}[closure of restricted Weil positivity]
\label{v4-arith:w5-direct:thm:restricted-RH}
Under the preceding pairing hypothesis, $q_W$ is nonnegative on
the closure of $D$ in $\mathscr A$. If that closure is all of
$\mathscr A$, RH follows. The same assertion holds in the
exponential space $\mathscr G$ of
\Cref{v4-arith:sc:gaussian-extension}.
\end{theorem}
\begin{proof}
Convolution, involution, and $W$ are continuous on these spaces.
Thus $f\mapsto q_W(f,f)$ is continuous and its nonnegative locus
is closed. Apply Weil's criterion when the closure is the full space.
\end{proof}

\begin{remark}[domain and spectrum]
\label{v4-arith:w5-direct:rem:what-not-proved}
The real even moment class $\mathscr A_{13}$ of
\Cref{v4-arith:pres57:moment-class} is a proper closed subspace
of the real even compact tests. Increasing a modular weight does
not remove its moment equations. An extension through Maass,
residual, or Eisenstein sectors requires maps into the specified
spectral summands and the same full pairing equality. A separate
comparison is required for each sector. None is logically necessary
for an independent proof of the full Weil inequality.
\end{remark}
\section{Restricted-Class Comparison}
\label{v4-arith:w5-direct:construction~5-map}

The restricted-class comparison partitions the conditions of
\Cref{v4-arith:w5-direct:rem:what-not-proved} into related
mathematical comparisons.

\begin{center}
\small
\begin{tabular}{p{0.45\textwidth}|p{0.50\textwidth}}
Condition & Reduction \\
\hline
F1.a.1 Banerjee--Clozel ideal &
\Cref{v4-arith:w5-e:chapter} \\
F1.a.2 small-prime extension (\(p=2,3\)) &
\Cref{v4-arith:w5-e:chapter} \\
F1.a.3.A block-Hochschild comparison &
\Cref{v4-arith:w5-f:F1a3-AB} \\
F1.a.3.B Paskunas \(\mathrm{Ext}^{1}\) vanishing &
\Cref{v4-arith:w5-f:F1a3-AB} \\
F1.a.3.C wild-EG compact generation (\(p=2,3\)) &
\Cref{v4-arith:w5-g:f1a3-C} \\
\(\mathrm{L}_{\mathrm{ACD}}\) adelic categorical descent &
\Cref{v4-arith:w5-e:chapter} \\
BC-diamond-glue (Ch1) Chenevier pseudochar &
\Cref{v4-arith:w5-h:bcglue-core-and-F2c} \\
BC-diamond-glue (AG.det) Arinkin--Gaitsgory split &
\Cref{v4-arith:w5-h:bcglue-core-and-F2c} \\
F2.c restricted-tensor \(E_{2}\) transport &
\Cref{v4-arith:w5-h:bcglue-core-and-F2c} \\
Arith.~Positselski: continuous duals &
\Cref{v4-arith:w5-d:chapter} \\
Arith.~Positselski: derived completion &
\Cref{v4-arith:w5-d:chapter} \\
Arith.~Positselski: detection on complete cofibers &
\Cref{v4-arith:w5-d:chapter} \\
Arith.~Positselski: Koszul-dual BC &
\Cref{v4-arith:w5-d:chapter} \\
Full \(\mathcal{V}^{\mathrm{prim}}\) P3: non-tempered (Ramanujan) &
\Cref{v4-arith:w5-c:full-vprim-P3} \\
Full \(\mathcal{V}^{\mathrm{prim}}\) P3: Eisenstein density correction &
\Cref{v4-arith:w5-c:full-vprim-P3} \\
H-P\(^{\mathrm{Ar}}\) (W) weight existence &
\Cref{v4-arith:w5-b:hp-operator} \\
(H) Weil formula at Gaussians &
\Cref{v4-arith:w5-i:chapter} \\
R\(_{\mathrm{analytic}}\) density-of-zeros frontier &
\Cref{v4-arith:w5-i:chapter} \\
R1.formal.b/c formal-Noetherian residuals &
\Cref{v4-arith:w5-i:chapter} \\
VB\(^{\ast}\) main direction primary Fock &
\Cref{v4-arith:w5-j:VBstar-spectral-positivity} \\
Restricted Weil pairing realization &
\Cref{v4-arith:w5-direct:thm:WP-isometry} \\
\end{tabular}
\end{center}


\section{The full Weil cone}
\label{v4-arith:w5-direct:irreducible-residual}
\begin{theorem}[full-domain Weil positivity]
\label{v4-arith:w5-direct:thm:irreducible-residual}
RH is equivalent to
\begin{equation}
\label{v4-arith:w5-direct:eq:irreducible}
 W(f*f^*)\ge0\qquad(f\in\mathscr A).
\end{equation}
It is also equivalent to existence of a positive Hilbert realization
of this full form. A realization through a specified arithmetic
sector requires the native maps and pairing equality in addition.
\end{theorem}
\begin{proof}
Use Weil's criterion and
\Cref{v4-arith:w5-direct:thm:WP-isometry} with $D=\mathscr A$.
\end{proof}
\begin{remark}[domain of the reduction]
\label{v4-arith:w5-direct:rem:reduction-honesty}
An identity on a proper restricted class does not supply this
full-domain assertion. Density must use a topology in which
$P,W_\infty,W_{\mathrm{fin}}$ and convolution are continuous.
The compact and exponential topologies above have that property.
\end{remark}
\section{Classical boundary}
\label{v4-arith:w5-direct:recap}
The functional equation, Hochschild character, Koszul pairing,
Bost--Connes dynamics, and coefficient multipliers address different
comparison problems. The coefficient construction gives a finite
holomorphic series. The integral pairing is positive on its actual
$L^2$ carrier. Their arithmetic relation requires specified maps.
For a fixed restricted carrier, a proved pairing identity gives
Weil positivity there and on its closure. The full Weil inequality
is RH-equivalent. A full abstract Hilbert realization follows from
that inequality by completion, while a realization in the specified
native arithmetic sector retains its additional structural requirements.
